% Make to Innovate Milestone Template
% Modified from the UCT Project report by Linus C. Brendel
% https://www.overleaf.com/latex/templates/uct-report-template/grctkzjtrqrm

%------------------------------------------------------
% DO NOT MODIFY THE FOLLOWING
\documentclass[12pt]{article}
\usepackage[english]{babel}
\usepackage{natbib}
\bibliographystyle{chicago}
\usepackage{url}
\usepackage{listings}
\usepackage{color} %red, green, blue, yellow, cyan, magenta, black, white
\usepackage[utf8x]{inputenc}
\usepackage{amsmath}
\usepackage{graphicx}
\graphicspath{{images/}}
\usepackage{parskip}
\usepackage{fancyhdr}
\usepackage{vmargin}
\setmarginsrb{3 cm}{2.5 cm}{3 cm}{2.5 cm}{1 cm}{1.5 cm}{1 cm}{1.5 cm}
%---------------------------------------------------------------------
% If you are including code snips or putting code in your appendix, you can define
% the colors used here.  See the listings example in the appendix for including 
% source code such as MatLab or Arduino.
\definecolor{mygreen}{RGB}{28,172,0} % color values Red, Green, Blue
\definecolor{mylilas}{RGB}{170,55,241}

%---------------------------------------------------------------------
% MODIFY THE FOLLOWING ITEMS FOR YOUR TITLE PAGE
%---------------------------------------------------------------------

% This is your title for your report.  It should include the milestone number and what this Milestone is about.
\title{M:2:I Milestone Example Report}
%This should be the team leader
\author{Matthew E. Nelson}
% Here enter all the names that contributed to this report.  This should be all the team members.
\newcommand{\members}{James Benson \ \\ Christine Nelson}
%This will automatically put today's date
\date{\today}

%Here is the role for each person listed.  It is in the same order as the author and then members.
\newcommand{\role}{Team Leader \ \\ Team Member \ \\ Team Member}
%Insert your faculty adviser here
\newcommand{\faculty}{Matthew E. Nelson}
%Insert your project name here
\newcommand{\project}{HABET}
%Iinsert your team name here
\newcommand{\team}{Engineering Team}

%------------------------------------------------------------------------
%DO NOT MODIFY ANY OF THE FOLLOWING ITEMS
%------------------------------------------------------------------------
\makeatletter
\let\thetitle\@title
\let\theauthor\@author
\let\thedate\@date
\makeatother

\pagestyle{fancy}
\fancyhf{}
\rhead{\theauthor}
\lhead{\thetitle}
\cfoot{\thepage}

\begin{document}

\begin{titlepage}
	\centering
    \includegraphics[scale = 0.4]{M2I_Boeing_Logo.png}\\[1.0 cm]
    \textsc{\LARGE Iowa State University}\\[1.0 cm]	
	\textsc{\Large AerE 294X/AerE 494X}\\[0.5 cm]
	\textsc{\large Make to Innovate}\\[0.5 cm]
	\rule{\linewidth}{0.2 mm} \\[0.4 cm]
	{ \huge \bfseries \thetitle}\\
	\rule{\linewidth}{0.2 mm} \\[1.5 cm]
    \emph{Project:}
    \project \ \\
    \emph{Team:}
    \team
	
	\begin{minipage}{0.4\textwidth}
		\begin{flushleft} \large
			\emph{Author(s):}\\
			\theauthor \ \\
            \members
			\end{flushleft}
			\end{minipage}~
			\begin{minipage}{0.4\textwidth}
			\begin{flushright} \large
			\emph{Role:} \\
			\role
		\end{flushright}
	\end{minipage}\\[1 cm]
	{\large Faculty Adviser: \faculty}\\[1 cm]
	{\large \thedate}\\[1 cm]
 
	\vfill
	
\end{titlepage}

%%%%%%%%%%%%%%%%%%%%%%%%%%%%%%%%%%%%%%%%%%%%%%%%%%%%%%%%%%%%%%%%%%%%%%%%%%%%%%%%%%%%%%%%

\tableofcontents
\pagebreak

%%%%%%%%%%%%%%%%%%%%%%%%%%%%%%%%%%%%%%%%%%%%%%%%%%%%%%%%%%%%%%%%%%%%%%%%%%%%%%%%%%%%%%%%

%-----------------------------------------------------------------------
% BEGIN YOUR DOCUMENT HERE
%-----------------------------------------------------------------------

%-----------------------------------------------------------------------
% INSTRUCTIONS
% This is where you will write your document.  Your table of contents will automatically get filled out based what you label as a section or subsection.  In addition, any paper, book, article or web site that you use must be cited.  This can be done using the biblist.bib file.  I have included all sections that are required and some examples.  Any images you use should go in the images folder.  This file will automatically search for images in that folder.
%-----------------------------------------------------------------------

\section{How to use this template}

This template is provided so that M:2:I teams can easily write reports that comply with the reporting standards set forth by the Make to Innovate Program Coordinator.  This template uses \LaTeX to generate a PDF that can be submitted on Blackboard.  Students should also upload their \LaTeX files as a zip file or should share them with the instructor.  This template will be made available through Overleaf as a template.

To begin with, make sure you edit the items noted in the above code.  You will want to make sure you that you set the Title, Authors (team members) and faculty adviser as it pertains to your group.  The rest should be left alone.  All other formatting is handled automatically by \LaTeX. 

\emph{Do not forget to delete this entire section.  It is only here to help you get started.  When done, simply delete all the text in this section.  The other sections will be required in your report.}

\subsection{Writing Style}
For writing a report students should read and refer to the book "Writing Style and Standards in Undergraduate Reports".  Specifically students should look at section 1.5 (page 32) and the example found on section 1.9 (page 55) of \cite{Donnell}.  Students should also refer to Chapter 2 which goes into more detail on writing standards.  Remember that you need to convey to your reader (in this case the instructor) about what you have accomplished, why you did what you did and how you came to any conclusions.  Your reports should have a good flow and it is encouraged that you use paragraphs to transition from one section to the next.

\subsection{Table of Contents}

The table of contents will be automatically created based on the use of \texttt{\textbackslash section\{\ldots\}} and \texttt{\textbackslash subsection\{\ldots\}} commands used.  This template will provide section headers for any section that is required.  You are free to add any other sections or subsections as needed.

\subsection{Images}

Images are handled differently in \LaTeX.  To begin with, you should upload all images into the "images" folder (if you are using Overleaf, otherwise place them in the images folder).  This file is setup so that it will search that folder for images.  To use images we first use the \texttt{\textbackslash begin\{figure\}} to setup an area for our image.  The \texttt{\textbackslash includegraphics[width=\textbackslash textwidth]\{sdr\_soilmoisture.png\}} is then used to call the image.  The "width" in the first part is set as the same width as your document text.  The filename is then next.  Finally we want to put a caption for this image.  This is done using \texttt{\textbackslash caption\{\ldots\}}.  It will automatically number the figure for you.  Finally, please note that \texttt{\textbackslash label\{\ldots\}} command.  This allows you to reference that image easily by using \texttt{\textbackslash ref\{\ldots\}} like this: Figure \ref{fig:sdr_graph} shows the calibrated soil moisture readings in comparison to the noise temperature.

\begin{figure}
\includegraphics[width=\textwidth]{sdr_soilmoisture.png}
\caption{An example image placed in the report.}
\label{fig:sdr_graph}
\end{figure}

A final note on images.  \LaTeX will automatically place the image in the best location according to the standards we have put forth.  This can be confusing as it may not put it exactly where you put the code.  However this perfectly acceptable and this is the reason why you need to reference your figures in your writing.  Generally \LaTeX will place images on the top of the page but it may move it to its own page if needed.  In other words, do not worry about the image placement, just let \LaTeX do its job.

\subsection{Bibliography}

Students are expected to cite all references used.  This includes any equations you used from any textbook, any journal articles \citep{Mohan2010ProfessionalCurriculum}, conference proceedings \citep{Nelson2014TeachingChipKits} websites \citep{ABET2015CriteriaPrograms} or other source material.  Make to Innovate uses the Chicago style for references.  The formatting will be handled for you automatically by \LaTeX.  However, you will need to add your references to the \texttt{biblist.bib} file.  You can upload and use other sources such as a link to your Mendeley account if you are using Overleaf.  Just make sure you edit the \texttt{\textbackslash bibliography\{\ldots\}} command that is at the bottom.

\subsection{Appendices}
Students are encouraged to use appendices in order to show additional information that is relative to the report.  This may include source code, additional 3D drawings, specifications and requirements, datasheets, or other additional information. 

\subsection{Final advice}
Remember that your report needs to tell a story, a story about the work that you did.  Plan your report out ahead of time and it should go a lot easier than just writing things out.  Also, yes, you will repeat yourself...a lot.  This is normal for almost all report writing.  Just make sure you are not just copying and pasting whole paragraphs, rephrase and rewrite.  The sections provided should give you a good start.  Make sure you do a spell check and also read it and have others read it for grammar mistakes.  Overleaf does have a spell checker, but does not check for grammar.  Finally if you want me (the instructor) to look over a report before submission I will be more than happy to do so.  Just remember that I will probably need at least 1 or 2 days to read over it and give you feedback.  Make sure you send me the source .tex file or share it with me on Overleaf.  Good luck!

%-------------------------------------------------------------------------
% Everything past here is required.  Do not delete these sections
% You can add to it if you wish.
%-------------------------------------------------------------------------
\newpage
\section*{Abstract}
\addcontentsline{toc}{section}{Abstract}

This is your abstract.  A well-structured abstract defines the need, specific goal or problem, the action you took, and the results and assessment of the results.  Your abstract should include what the deliverable(s) of the milestone was, what the milestone is about and why it is important.  The abstract should not be numbered and should be on its own page.

\newpage
\section{Introduction}
This is your introduction.  It should state the problem and what the need is in general terms.  Introduce your problem and how this milestone will solve it.

\section{Background}
A background of the problem or system.  Describe what lead to this issue if that applies.  Your goal should also be clearly specified here.

\subsection{Deliverable}
Outline what exactly the deliverable(s) are for this milestone.

\newpage
\section{Problem Identification}
Fully outline your problem statement.  Include all technical details in this section that helps to define the problem that this milestone will solve.

\section{Problem Solution}
State your solution to the problem.  Include all technical information here.

% Most teams will probably have some sort of experiment to verify that their system is working properly.  However, if you do not, you still need to report on how you verified that your deliverable meet the requirements.  You can change these headings if you need to, but you should have something here.

\section{Experimental Validation}
Your deliverable should include that you have built and/or tested something.  Include those results here.

\subsection{Experimental Theory}
What were your expected results?  State what theory backups your claim.


\subsection{Experimental Procedure}
This should describe what procedures you did to conduct your test or experiment.

\subsection{Experimental Results}
This is the results of that testing or experimentation.

\newpage
% This is the team's opportunity to tell me more about the milestone, how things went, what they learned, etc.  
\section{Discussion}
Here discuss how successful the milestone and deliverable are and other thoughts that you have in regards to this milestone.  This is also a good place to make recommendations for future work if needed.

\section{Conclusion}

Summarize your report here.

\section{Acknowledgments}
If you wish to acknowledge any person, group or entity you should do so here.

%-------------------------------------------------------------------------
% Do not edit this unless you are using a different Bib file.  In that case
% put the bib file name below (do not include the .bib extension
%-------------------------------------------------------------------------

\newpage

\bibliography{biblist}

%-------------------------------------------------------------------------
% You may delete this if you are not using an appendix in your report.
%-------------------------------------------------------------------------

\newpage
\appendix
\section{Example Appendix}
This is a very short example of doing an appendix in \LaTeX.  An appendix is not required but is a good place to put code, additional drawings, tables or other information that will help the instructor in determining that you met the requirements of the milestone.  I have included an example of how you can use \LaTeX to automatically import code.

% The following tells listings how to treat the code.  You can edit this
% as needed

\lstset{language=Matlab,%
    %basicstyle=\color{red},
    breaklines=true,%
    morekeywords={matlab2tikz},
    keywordstyle=\color{blue},%
    morekeywords=[2]{1}, keywordstyle=[2]{\color{black}},
    identifierstyle=\color{black},%
    stringstyle=\color{mylilas},
    commentstyle=\color{mygreen},%
    showstringspaces=false,%without this there will be a symbol in the places where there is a space
    numbers=left,%
    numberstyle={\tiny \color{black}},% size of the numbers
    numbersep=9pt, % this defines how far the numbers are from the text
    emph=[1]{for,end,break},emphstyle=[1]\color{red}, %some words to emphasise
    %emph=[2]{word1,word2}, emphstyle=[2]{style},    
}

% To import your code, just use the following command.  Make sure it is the same
% folder or if using Overleaf it is uploaded to the project.

\lstinputlisting{lab1_hello.m}

%-------------------------------------------------------------------------
% DO NOT DELETE THIS.
% This is the end of the document.  Nothing else should go past this.
%-------------------------------------------------------------------------

\end{document}